% Chapter 6: verification

\chapter{Verification}

\label{Chapter6}

%----------------------------------------------------------------------------------------

\section{Lattice identities checked by the tests}

\file{tests/unit/test\_lattices.py} is parametrised over every implemented lattice and over
fp32 and fp64. The tests assert
\begin{align}
\langle 1 \rangle &= 1, \\
\langle c_\alpha \rangle = \langle c_\alpha c_\beta c_\gamma \rangle &= 0, \\
\langle c_\alpha c_\beta \rangle &= \cs^2\,\delta_{\alpha\beta}, \\
\langle c_\alpha c_\beta c_\gamma c_\delta \rangle &= \cs^4\,\Delta^{(4)}_{\alpha\beta\gamma\delta}
  &&\text{if } \code{isotropy\_order} \ge 4, \\
\langle c_{\alpha_1}\cdots c_{\alpha_5} \rangle = 0,\quad
\langle c_{\alpha_1}\cdots c_{\alpha_6} \rangle &= \cs^6\,\Delta^{(6)}
  &&\text{if } \code{isotropy\_order} \ge 6, \\
\langle c_{\alpha_1}\cdots c_{\alpha_7} \rangle = 0,\quad
\langle c_{\alpha_1}\cdots c_{\alpha_8} \rangle &= \cs^8\,\Delta^{(8)}
  &&\text{if } \code{isotropy\_order} \ge 8,
\end{align}
where $\Delta^{(2n)}$ is the fully symmetric isotropic tensor, the sum over all pairings of
the $2n$ indices of products of Kronecker deltas (3 terms at fourth order, 15 at sixth, 105
at eighth). The test builds $\Delta^{(2n)}$ generically by enumerating pairings. Tolerances
are in \cref{tab:tolerances}.

\begin{table}[h]
\centering\small
\begin{tabular}{@{}lcc@{}}
\toprule
Identity & fp32 tolerance & fp64 tolerance \\
\midrule
orders 0--3 & $10^{-6}$ & $10^{-12}$ \\
order 4     & $10^{-5}$ & $10^{-11}$ \\
orders 5--6 & $10^{-4}$ & $10^{-10}$ \\
orders 7--8 & $10^{-3}$ & $10^{-9}$ \\
\bottomrule
\end{tabular}
\caption{Absolute tolerances. Higher orders involve $|\vc|^{8}$ up to $3^8$ for D2Q37, hence
the looser bounds.}
\label{tab:tolerances}
\end{table}

Further checks: \code{opp} is an int64 involution with $\vc_{\bar i} = -\vc_i$ and
$w_{\bar i} = w_i$; velocities are distinct integers with the rest velocity first;
\code{shifts} equals \code{c} as Python ints; constants are built in the requested dtype;
fp64 constants differ from upcast fp32 ones; \code{.to(float64)} casts only the float
buffers; the state dict is empty and there are no parameters; the module moves to the device
\code{get\_device()} selects; D2Q5 fails fourth-order isotropy; D2Q37's second moment equals
$1/r^2$; non-integer velocities and mismatched shapes raise \code{ValueError}. One CUDA-only
test checks \code{f[opp]} indexing on the GPU.

%----------------------------------------------------------------------------------------

\section{Moments and equilibrium checked by the tests}

\file{tests/unit/test\_equilibrium.py} is parametrised over every implemented lattice and
over fp32 and fp64, on a $4\times5$ (2D) or $3\times4\times5$ (3D) grid with a seeded
density $\rho \in [1, 1.1]$ and velocity $|u_\alpha| \le 0.05$. Absolute tolerances are
$10^{-5}$ in fp32 and $10^{-12}$ in fp64.

For \code{moments()}: output shapes and dtype; rest populations $w_i\rho_0$ return
$\rho_0$ and $\vu = 0$; and fp32 moments of an fp64 field rounded to fp32 agree with the
fp64 reference to a relative $2\times10^{-7}$ in $\rho$ and $10^{-6}$ in $\vu$, which is
fp32 rounding of the result rather than fp32 summation error.

For \code{equilibrium()} at the lattice's default order, with
$\langle\cdot\rangle_{\mathrm{eq}} = \sum_i (\cdot)\, f^{\mathrm{eq}}_i$:
\begin{align}
\langle 1 \rangle_{\mathrm{eq}} &= \rho, \qquad
\langle c_\alpha \rangle_{\mathrm{eq}} = \rho u_\alpha, \\
\langle c_\alpha c_\beta \rangle_{\mathrm{eq}}
  &= \rho\,(\cs^2\delta_{\alpha\beta} + u_\alpha u_\beta)
  &&\text{if } \code{isotropy\_order} \ge 4, \\
\langle c_\alpha c_\beta c_\gamma \rangle_{\mathrm{eq}}
  &= \rho\,u_\alpha u_\beta u_\gamma
     + \rho\cs^2\,(u_\alpha\delta_{\beta\gamma} + \text{2 perms})
  &&\text{if } \code{equilibrium\_order} \ge 3, \\
\langle c_\alpha c_\beta c_\gamma c_\delta \rangle_{\mathrm{eq}}
  &= \rho\,u_\alpha u_\beta u_\gamma u_\delta
     + \rho\cs^2\,(u_\alpha u_\beta\delta_{\gamma\delta} + \text{5 perms}) \nonumber \\
  &\qquad + \rho\cs^4\Delta^{(4)}_{\alpha\beta\gamma\delta}
  &&\text{if } \code{equilibrium\_order} \ge 4,
\end{align}
the Maxwell--Boltzmann moments at the reference temperature. The second-moment identity is
skipped on D2Q5 by design (\cref{Chapter3}); the third and fourth run only on D2Q37.
Further checks: output shape and dtype; $f^{\mathrm{eq}} = w_i\rho$ at rest;
\code{moments(equilibrium(rho, u))} returns \code{(rho, u)}; $f^{\mathrm{eq}} > 0$ at this
Mach number; the default order equals \code{lattice.equilibrium\_order}; order~0 and
\code{equilibrium\_order + 1} raise \code{ValueError}; and the output is on the device of
the inputs.

\file{tests/gradcheck/test\_equilibrium\_grad.py} runs in fp64 on D2Q9, D2Q37 and D3Q19
with $3\times4$ or $2\times3\times3$ grids. \code{torch.autograd.gradcheck} passes for
\code{moments()} in \code{f} and for \code{equilibrium()} in \code{(rho, u)} jointly, and a
backward pass through $\sum (f^{\mathrm{eq}})^2$ reaches both inputs with finite gradients.
A failure here means an in-place operation, a \code{.detach()} or a non-differentiable cast
has entered the path; it does not test the physics.

\file{tests/unit/test\_lattices.py} additionally asserts that \code{equilibrium\_order} is
an integer with $1 \le \code{equilibrium\_order}$ and
$2\,\code{equilibrium\_order} \le \code{isotropy\_order}$ on every lattice, and that a
subclass violating the bound raises \code{ValueError} at construction.

\section{Streaming checked by the tests}

\file{tests/unit/test\_streaming.py} is parametrised over every implemented lattice and over
fp32 and fp64, on a $5\times7$ (2D) or $4\times5\times6$ (3D) grid with a seeded positive
field $f \in [0.5, 1.5)$ of distinct values, so that a misplaced entry is visible. The grid
exceeds the longest shift on every lattice ($|c_\alpha| \le 3$ on D2Q37), so a wrap is a
genuine wrap and not a shift by a multiple of the grid. Because streaming is a permutation,
every comparison is exact (\code{torch.equal}); there are no tolerances.

For \code{stream()}:
\begin{itemize}[nosep]
\item output shape, dtype and device match the input; the output is a new tensor and the
      input is unchanged;
\item \textbf{push}: for every direction $q$, the value at the origin arrives at
      $\vc_q \bmod$ grid, and the whole slab equals \code{torch.roll(f[q], shifts[q])};
\item the single rest population is unchanged;
\item \textbf{periodic wrap}: for every direction and every axis with a positive shift, the
      value on the last node of that axis arrives on node $c_\alpha - 1$;
\item the sorted entries of the output equal the sorted entries of the input (a
      permutation); rolling every direction back by $-\vc_q$ recovers $f$ exactly; and
      rolling direction $q$ by $\vc_{\bar{q}}$ also recovers $f_q$, so the \code{opp} table
      and the shift list agree;
\item \textbf{conservation} (marker \code{conservation}, fp64): the per-direction sorted
      slabs match, hence the direction sums, hence the global mass and
      $\sum_q \vc_q \sum_{\mathbf{x}} f_q$ are exactly unchanged; and on a spatially uniform
      field streaming is the identity, so \code{moments()} returns identical $\rho$, $\vu$;
\item a wrong $Q$ or a wrong number of spatial axes raises \code{ValueError} naming the
      offending quantity; and, with CUDA, the result on the device equals the result on the
      host.
\end{itemize}
The global conservation check sorts before summing because summation order differs between
the two fields; a last-bit difference in a direct fp32 sum is not a streaming error.

\file{tests/gradcheck/test\_streaming\_grad.py} runs in fp64 on D2Q9, D2Q37 and D3Q19 with
$3\times4$ or $2\times3\times3$ grids. \code{torch.autograd.gradcheck} passes; the
vector--Jacobian product equals the upstream gradient streamed with every shift negated,
exactly, since the Jacobian of a permutation is a permutation matrix; and a five-step
composed rollout stays differentiable with $\partial \sum f'^2 / \partial f = 2f$. A failure
here means an in-place write or a non-differentiable op has entered the path.

%----------------------------------------------------------------------------------------

\section{Device utilities}

\file{tests/unit/test\_gpu.py} covers device resolution and the environment override, both
branches of \code{accumulation\_dtype}, the query functions, byte formatting, the memory
no-ops on CPU, \code{track\_memory} on CPU, and \code{Timer}. Three tests carry the
\code{gpu} marker and run only with CUDA: GPU info fields, peak-memory tracking, and the
CUDA device summary.

%----------------------------------------------------------------------------------------

\section{Current status}

\begin{table}[h]
\centering\small
\begin{tabular}{@{}lr@{}}
\toprule
\code{pytest} (whole repository) & 454 passed, 65 skipped \\
\code{ruff check}, \code{ruff format --check} & clean \\
\code{mypy fluxlb} & clean (40 files) \\
\bottomrule
\end{tabular}
\caption{Test and static-check status at the time of writing. Skips are the CUDA-only tests,
the higher-order isotropy identities on lattices that do not claim them, and the
higher-order equilibrium moments on lattices below that order.}
\end{table}

%----------------------------------------------------------------------------------------

\section{Adding a lattice}

\begin{enumerate}[nosep]
\item Subclass \code{Lattice}; set \code{Q}, \code{D}, \code{isotropy\_order} and
      \code{equilibrium\_order} $= \lfloor\code{isotropy\_order}/2\rfloor$.
\item In \code{\_\_init\_\_}, build \code{c} (int64), \code{w} (float64 from exact
      fractions) and \code{opp} (int64, preferably derived from the velocity list), then call
      \code{super().\_\_init\_\_(c=c, w=w, opp=opp, cs2=..., dtype=dtype)}.
\item Wrap hand-laid-out tables in \code{\# fmt: off} / \code{\# fmt: on} so the formatter
      keeps the shell layout.
\item State the ordering convention in the docstring; \code{opp} and the boundary code depend
      on it.
\item Add the class to \code{LATTICES} in the test file. The isotropy tests will reject a
      wrong weight or velocity immediately.
\end{enumerate}
